% GENERATED FILE. Edit canonical lessons, then run npm run generate:books.
% canonical-source-hash: fnv1a64:a3bc5d96b855300e
% canonical-lessons: KA-C405-maitri, KA-C405-sandhana, KA-C405-madhyasthike, KA-C405-hondanike, KA-C405-pracodisu

\chapter{Settling Disputes}
\label{ch:ka-settling-disputes}
\hlchaptermodality{405}

\begin{chapteropening}
\textbf{By the end of this chapter:} \emph{I can say an alliance, a negotiation, mediation, a compromise and to incite.}

Complete the last lesson of chapter 405: I can say an alliance, a negotiation, mediation, a compromise and to incite.
\end{chapteropening}

\section[\texorpdfstring{\kn{ಮೈತ್ರಿ} (maitri)}{ಮೈತ್ರಿ (maitri)}]{\kn{ಮೈತ್ರಿ} (maitri) — an alliance; friendship}

\label{lesson:KA-C405-maitri}



\begin{quote}
Before the new one: say the Kannada for to coordinate, then the Kannada for goodwill.
\end{quote}

\subsection*{You'll want to know: \kn{ಮೈತ್ರಿ}}
\textbf{\kn{ಮೈತ್ರಿ}} — \emph{maitri} — \textquotedblleft{}an alliance; friendship\textquotedblright{}.

\begin{quote}
\kn{ಎರಡು} \kn{ಊರುಗಳ} \kn{ನಡುವೆ} \kn{ಹಳೆಯ} \kn{ಮೈತ್ರಿ} \kn{ಇದೆ}.
\end{quote}

\begin{quote}
\emph{eraḍu ūrugaḷa naḍuve haḷeya maitri ide.}
\end{quote}

\begin{quote}
There is an old friendship between the two towns.
\end{quote}

Formal; between groups or countries.

\begin{cousinweb}
\textbf{\kn{ಮೈತ್ರಿ}} (\emph{maitri}) is borrowed from Sanskrit \emph{maitrī}, \textquotedblleft{}friendship, goodwill\textquotedblright{} (Monier-Williams), from \emph{mitra}, \textquotedblleft{}a friend\textquotedblright{}.
\end{cousinweb}

\begin{grammarlens}[title={What you've built}]
One more word for giving reasons and opinions.
\end{grammarlens}

\subsection*{Your turn}
\begin{itemize}
\raggedright
  \item \emph{Say it:} \emph{maitri}
  \item \emph{Say it:} \emph{maitri}, once more
  \item \emph{Say it:} \emph{eraḍu ūrugaḷa naḍuve haḷeya maitri ide.} — There is an old friendship between the two towns.
  \item \emph{Recall:} say the Kannada for to coordinate, then the Kannada for goodwill, then say \emph{maitri} again
\end{itemize}

\begin{tcolorbox}[breakable,colback=teal!4,colframe=teal!35!black,title={Before you move on}]
What does \kn{ಮೈತ್ರಿ} mean? (\textquotedblleft{}An alliance; friendship\textquotedblright{}.) Say it once more.
\end{tcolorbox}
\section[\texorpdfstring{\kn{ಸಂಧಾನ} (sandhāna)}{ಸಂಧಾನ (sandhāna)}]{\kn{ಸಂಧಾನ} (sandhāna) — a negotiation}

\label{lesson:KA-C405-sandhana}



\begin{quote}
Before the new one: say the Kannada for goodwill, then the Kannada for an alliance.
\end{quote}

\subsection*{You'll want to know: \kn{ಸಂಧಾನ}}
\textbf{\kn{ಸಂಧಾನ}} — \emph{sandhāna} — \textquotedblleft{}a negotiation\textquotedblright{}.

\begin{quote}
\kn{ಎರಡು} \kn{ಕುಟುಂಬಗಳ} \kn{ನಡುವೆ} \kn{ಸಂಧಾನ} \kn{ನಡೆಯಿತು}.
\end{quote}

\begin{quote}
\emph{eraḍu kuṭumbagaḷa naḍuve sandhāna naḍeyitu.}
\end{quote}

\begin{quote}
Talks took place between the two families.
\end{quote}

\kn{ಸಂಧಾನಕಾರ}, a negotiator.

\begin{cousinweb}
\textbf{\kn{ಸಂಧಾನ}} (\emph{sandhāna}) is borrowed from Sanskrit \emph{saṃdhāna}, \textquotedblleft{}joining, uniting\textquotedblright{} (Monier-Williams): bringing two sides together.
\end{cousinweb}

\begin{grammarlens}[title={What you've built}]
One more word for giving reasons and opinions.
\end{grammarlens}

\subsection*{Your turn}
\begin{itemize}
\raggedright
  \item \emph{Say it:} \emph{sandhāna}
  \item \emph{Say it:} \emph{sandhāna}, once more
  \item \emph{Say it:} \emph{eraḍu kuṭumbagaḷa naḍuve sandhāna naḍeyitu.} — Talks took place between the two families.
  \item \emph{Recall:} say the Kannada for goodwill, then the Kannada for an alliance, then say \emph{sandhāna} again
\end{itemize}

\begin{tcolorbox}[breakable,colback=teal!4,colframe=teal!35!black,title={Before you move on}]
What does \kn{ಸಂಧಾನ} mean? (\textquotedblleft{}A negotiation\textquotedblright{}.) Say it once more.
\end{tcolorbox}
\section[\texorpdfstring{\kn{ಮಧ್ಯಸ್ಥಿಕೆ} (madhyasthike)}{ಮಧ್ಯಸ್ಥಿಕೆ (madhyasthike)}]{\kn{ಮಧ್ಯಸ್ಥಿಕೆ} (madhyasthike) — mediation}

\label{lesson:KA-C405-madhyasthike}



\begin{quote}
Before the new one: say the Kannada for an alliance, then the Kannada for a negotiation.
\end{quote}

\subsection*{You'll want to know: \kn{ಮಧ್ಯಸ್ಥಿಕೆ}}
\textbf{\kn{ಮಧ್ಯಸ್ಥಿಕೆ}} — \emph{madhyasthike} — \textquotedblleft{}mediation\textquotedblright{}.

\begin{quote}
\kn{ಊರಿನ} \kn{ಹಿರಿಯರ} \kn{ಮಧ್ಯಸ್ಥಿಕೆಯಿಂದ} \kn{ಜಗಳ} \kn{ಮುಗಿಯಿತು}.
\end{quote}

\begin{quote}
\emph{ūrina hiriyara madhyasthikeyinda jagaḷa mugiyitu.}
\end{quote}

\begin{quote}
The quarrel ended through the village elders' mediation.
\end{quote}

\kn{ಮಧ್ಯಸ್ಥ}, a mediator.

\begin{cousinweb}
\textbf{\kn{ಮಧ್ಯಸ್ಥಿಕೆ}} (\emph{madhyasthike}) is built on Sanskrit \emph{madhyastha}, \textquotedblleft{}being in the middle\textquotedblright{} (Monier-Williams), with the Kannada noun ending -\emph{ike}.
\end{cousinweb}

\begin{grammarlens}[title={What you've built}]
One more word for giving reasons and opinions.
\end{grammarlens}

\subsection*{Your turn}
\begin{itemize}
\raggedright
  \item \emph{Say it:} \emph{madhyasthike}
  \item \emph{Say it:} \emph{madhyasthike}, once more
  \item \emph{Say it:} \emph{ūrina hiriyara madhyasthikeyinda jagaḷa mugiyitu.} — The quarrel ended through the village elders' mediation.
  \item \emph{Recall:} say the Kannada for an alliance, then the Kannada for a negotiation, then say \emph{madhyasthike} again
\end{itemize}

\begin{tcolorbox}[breakable,colback=teal!4,colframe=teal!35!black,title={Before you move on}]
What does \kn{ಮಧ್ಯಸ್ಥಿಕೆ} mean? (\textquotedblleft{}Mediation\textquotedblright{}.) Say it once more.
\end{tcolorbox}
\section[\texorpdfstring{\kn{ಹೊಂದಾಣಿಕೆ} (hondāṇike)}{ಹೊಂದಾಣಿಕೆ (hondāṇike)}]{\kn{ಹೊಂದಾಣಿಕೆ} (hondāṇike) — a compromise; an adjustment}

\label{lesson:KA-C405-hondanike}



\begin{quote}
Before the new one: say the Kannada for a negotiation, then the Kannada for mediation.
\end{quote}

\subsection*{You'll want to know: \kn{ಹೊಂದಾಣಿಕೆ}}
\textbf{\kn{ಹೊಂದಾಣಿಕೆ}} — \emph{hondāṇike} — \textquotedblleft{}a compromise; an adjustment\textquotedblright{}.

\begin{quote}
\kn{ಸ್ವಲ್ಪ} \kn{ಹೊಂದಾಣಿಕೆ} \kn{ಮಾಡಿಕೊಂಡರೆ} \kn{ಸಮಸ್ಯೆ} \kn{ಮುಗಿಯುತ್ತದೆ}.
\end{quote}

\begin{quote}
\emph{svalpa hondāṇike māḍikoṇḍare samasye mugiyuttade.}
\end{quote}

\begin{quote}
If we compromise a little, the problem ends.
\end{quote}

\kn{ಹೊಂದಾಣಿಕೆ} \kn{ಮಾಡಿಕೊ}, to adjust, to compromise.

\begin{cousinweb}
\textbf{\kn{ಹೊಂದಾಣಿಕೆ}} (\emph{hondāṇike}) is native Dravidian, built on \emph{hondu}, \textquotedblleft{}to come together, to suit\textquotedblright{} (DEDR 4541).
\end{cousinweb}

\begin{grammarlens}[title={What you've built}]
One more word for giving reasons and opinions.
\end{grammarlens}

\subsection*{Your turn}
\begin{itemize}
\raggedright
  \item \emph{Say it:} \emph{hondāṇike}
  \item \emph{Say it:} \emph{hondāṇike}, once more
  \item \emph{Say it:} \emph{svalpa hondāṇike māḍikoṇḍare samasye mugiyuttade.} — If we compromise a little, the problem ends.
  \item \emph{Recall:} say the Kannada for a negotiation, then the Kannada for mediation, then say \emph{hondāṇike} again
\end{itemize}

\begin{tcolorbox}[breakable,colback=teal!4,colframe=teal!35!black,title={Before you move on}]
What does \kn{ಹೊಂದಾಣಿಕೆ} mean? (\textquotedblleft{}A compromise; an adjustment\textquotedblright{}.) Say it once more.
\end{tcolorbox}
\section[\texorpdfstring{\kn{ಪ್ರಚೋದಿಸು} (pracōdisu)}{ಪ್ರಚೋದಿಸು (pracōdisu)}]{\kn{ಪ್ರಚೋದಿಸು} (pracōdisu) — to incite}

\label{lesson:KA-C405-pracodisu}



\begin{quote}
Before the new one: say the Kannada for mediation, then the Kannada for a compromise.
\end{quote}

\subsection*{You'll want to know: \kn{ಪ್ರಚೋದಿಸು}}
\textbf{\kn{ಪ್ರಚೋದಿಸು}} — \emph{pracōdisu} — \textquotedblleft{}to incite\textquotedblright{}.

\begin{quote}
\kn{ಯಾರೋ} \kn{ಜನರನ್ನು} \kn{ಜಗಳಕ್ಕೆ} \kn{ಪ್ರಚೋದಿಸಿದರು}.
\end{quote}

\begin{quote}
\emph{yārō janarannu jagaḷakke pracōdisidaru.}
\end{quote}

\begin{quote}
Someone incited the people to fight.
\end{quote}

Past (I …): \kn{ಪ್ರಚೋದಿಸಿದೆ}, \emph{pracōdiside}.

The person takes -\kn{ಅನ್ನು}, the aim -\kn{ಗೆ}. \kn{ಪ್ರಚೋದನೆ}, provocation.

\begin{cousinweb}
\textbf{\kn{ಪ್ರಚೋದಿಸು}} (\emph{pracōdisu}) is built on Sanskrit \emph{pracodana}, \textquotedblleft{}instigating, urging\textquotedblright{} (Monier-Williams), with the Kannada verb ending -\emph{isu}.
\end{cousinweb}

\begin{grammarlens}[title={What you've built}]
One more word for giving reasons and opinions.
\end{grammarlens}

\subsection*{Your turn}
\begin{itemize}
\raggedright
  \item \emph{Say it:} \emph{pracōdisu}
  \item \emph{Say it:} \emph{pracōdisu}, once more
  \item \emph{Say it:} \emph{yārō janarannu jagaḷakke pracōdisidaru.} — Someone incited the people to fight.
  \item \emph{Recall:} say the Kannada for mediation, then the Kannada for a compromise, then say \emph{pracōdisu} again
\end{itemize}

\begin{tcolorbox}[breakable,colback=teal!4,colframe=teal!35!black,title={Before you move on}]
What does \kn{ಪ್ರಚೋದಿಸು} mean? (\textquotedblleft{}To incite\textquotedblright{}.) Say it once more.
\end{tcolorbox}
